\documentclass[a4paper,11pt]{article}
\usepackage[utf8]{inputenc}
\usepackage[T1]{fontenc}
\usepackage[ngerman]{babel}
\usepackage{amsmath,amssymb}
\usepackage{xcolor}
\usepackage{tikz}
\usepackage{array}
\usepackage[a4paper,margin=2.2cm]{geometry}

\definecolor{bgdark}{HTML}{1E1E1E}
\definecolor{fgtext}{HTML}{E6E6E6}
\definecolor{gridcol}{HTML}{4A4A4A}
\definecolor{accent}{HTML}{6CB6FF}
\definecolor{leicht}{HTML}{7BD88F}
\definecolor{mittel}{HTML}{E6C07B}
\definecolor{schwer}{HTML}{E06C75}
\pagecolor{bgdark}
\color{fgtext}

\newcommand{\karo}[1]{\par\noindent
  \begin{tikzpicture}
    \draw[step=5mm, gridcol, very thin] (0,0) grid (\linewidth,#1);
  \end{tikzpicture}\par}
\newcommand{\aufgabe}[4]{\vspace{0.6em}\par\noindent
  \textbf{\large Aufgabe #1}\quad #2 \hfill {\color{#3}\textbf{[#4]}}\par\smallskip}
\setlength{\parindent}{0pt}

\begin{document}

\begin{center}
  {\color{accent}\Huge\bfseries Analysis II}\\[0.3em]
  {\Large Übungsaufgaben 04 — Trigonometrie, Komplexe Zahlen \& Fourieranalyse}\\[0.8em]
\end{center}

\noindent\textbf{Name:} \rule{6cm}{0.4pt} \hfill \textbf{Datum:} \rule{3.5cm}{0.4pt}
\vspace{0.4em}
{\color{gridcol}\hrule height 1pt}
\vspace{0.3em}

\small\textit{Nützlich: $\sin^2 x+\cos^2 x=1$, $\tan x=\tfrac{\sin x}{\cos x}$, $\sec x=\tfrac{1}{\cos x}$;
Eulersche Formel $z=a+ib=r(\cos\varphi+i\sin\varphi)=r\,e^{i\varphi}$ mit
$r=\sqrt{a^2+b^2}$, $\varphi=\arctan(b/a)$.}
\normalsize
\vspace{0.4em}

{\color{accent}\large\bfseries Teil A — Trigonometrie}\par\vspace{0.3em}

%==================================================================
\aufgabe{1}{Bogenmaß}{leicht}{leicht}
Rechne in das jeweils andere Maß um:
$\ 60^\circ,\ 120^\circ,\ 270^\circ$ in Bogenmaß sowie
$\ \tfrac{\pi}{6}$ und $\tfrac{3\pi}{2}$ in Grad.
\karo{5cm}

%==================================================================
\aufgabe{2}{Wellenparameter}{mittel}{mittel}
Gegeben sei die Schwingung $f(t) = 3\sin(4\pi t)$. Bestimme Amplitude $A$,
Kreisfrequenz $\omega$, Frequenz $f$ und Wellenlänge $\lambda$.
\karo{5cm}

%==================================================================
\aufgabe{3}{Ableitungen trigonometrischer Funktionen}{mittel}{mittel}
Berechne die Ableitungen:
\[
  f(x)=\sin(x)\cos(x), \quad
  g(x)=\arctan(x), \quad
  h(x)=x\sin(x), \quad
  k(x)=\tanh(x).
\]
\karo{6.5cm}

%==================================================================
\aufgabe{4}{Trigonometrische Identitäten}{schwer}{schwer}
Zeige durch geeignete Umformungen:
\[
  \text{(a)}\quad 1 + \tan^2(x) = \sec^2(x)
  \qquad\qquad
  \text{(b)}\quad \cosh^2(x) - \sinh^2(x) = 1.
\]
\textit{Hinweis zu (b): $\sinh(x)=\tfrac{e^x-e^{-x}}{2}$, $\cosh(x)=\tfrac{e^x+e^{-x}}{2}$.}
\karo{6cm}

\vspace{0.4em}
{\color{accent}\large\bfseries Teil B — Komplexe Zahlen}\par\vspace{0.3em}
\textit{In den Aufgaben 5 und 6 sei $z_1 = 3+2i$ und $z_2 = 1-4i$.}
\vspace{0.3em}

%==================================================================
\aufgabe{5}{Grundrechenarten}{leicht}{leicht}
Berechne $z_1+z_2$, $z_1-z_2$ und $z_1\cdot z_2$. Gib außerdem $\operatorname{Re}(z_1)$,
$\operatorname{Im}(z_1)$ und die komplex Konjugierte $\overline{z_1}$ an.
\karo{6cm}

%==================================================================
\aufgabe{6}{Betrag und Division}{mittel}{mittel}
Berechne $|z_1|$ sowie den Quotienten $\dfrac{z_1}{z_2}$ (in der Form $a+ib$).
\karo{6cm}

%==================================================================
\aufgabe{7}{Polarform / Eulersche Formel}{mittel}{mittel}
Schreibe $z=1+i$ und $z=2i$ in Polarform $r\,e^{i\varphi}$. Wandle umgekehrt
$2\,e^{i\pi/3}$ in die Form $a+ib$ um.
\karo{6cm}

%==================================================================
\aufgabe{8}{Potenzieren mit der Eulerschen Formel}{schwer}{schwer}
Berechne $(1+i)^4$ mithilfe der Polarform und überprüfe das Ergebnis durch
direktes Ausmultiplizieren.
\karo{6cm}

\vspace{0.4em}
{\color{accent}\large\bfseries Teil C — Fourieranalyse}\par\vspace{0.3em}

%==================================================================
\aufgabe{9}{Konzeptfragen}{leicht}{leicht}
(a) Schreibe $A\cos(\omega t) + i\,A\sin(\omega t)$ als komplexe Exponentialfunktion. \\
(b) Was bildet die Fouriertransformation $F(\omega)=\int_{-\infty}^{\infty} f(t)\,e^{-i\omega t}\,dt$
    ineinander ab, und was beschreibt $|F(\omega)|$? \\
(c) Erkläre kurz konstruktive und destruktive Interferenz.
\karo{6cm}

%==================================================================
\aufgabe{10}{Werte der komplexen Exponentialfunktion}{mittel}{mittel}
Berechne $e^{i\pi}$, $e^{i\pi/2}$ und $e^{i2\pi}$ mithilfe der Eulerschen Formel
und gib die berühmte Identität für $e^{i\pi}$ an.
\karo{5.5cm}

\end{document}
